\section{Erd\H{o}s Problem \#252}

\subsection*{1) ROUND-2 OBJECTIVE}
Path (C): keep the successful Round-1 factorial-tail framework, but sharpen it into an exact finite-dimensional reduction for the unresolved cases.

Round-1 already gave complete proofs for $k=1$ and $k=2$. The genuine remaining issue is the general case, where the naive prime-tail estimate is too crude. The most promising Round-2 move is therefore to determine exactly how many local terms one must keep, and exactly how small the remaining tail becomes.

\subsection*{2) Round-1 FOUNDATION USED}
I use the following Round-1 results without re-proving them.

\begin{itemize}
\item If
\[
\alpha_k:=\sum_{n=1}^{\infty}\frac{\sigma_k(n)}{n!}
\]
is rational and $N$ is at least the denominator, then the factorial tail
\[
R_{N,k}:=\sum_{n=N+1}^{\infty}\sigma_k(n)\frac{N!}{n!}
\]
is an integer.
\item The cases $k=1$ and $k=2$ were completely settled in Round-1.
\end{itemize}

In this round I work only with $k\ge 2$ and with the prime choice $N=p-1$.

\subsection*{3) NEW INSIGHT / TOOL (ROUND-2)}
The new tool is an exact truncation threshold for the prime-tail method.

For a prime $p$ define
\[
R_{p-1,k}=\sum_{j=0}^{\infty}U_{j,k}(p),
\qquad
U_{j,k}(p):=\frac{\sigma_k(p+j)}{p(p+1)\cdots(p+j)}.
\]
Then:

\begin{itemize}
\item If one truncates after fewer than $k$ local terms, the omitted tail is still at least $1$ (in fact usually much larger).
\item If one keeps exactly the first $k$ local terms, the remaining tail is already $O_k(1/p)$.
\end{itemize}

Thus the irrationality problem for fixed $k\ge 2$ is reduced to understanding the near-integrality of the single finite expression
\[
Q_{p,k}:=\sum_{j=0}^{k-1}U_{j,k}(p)
\]
along large primes $p$.

\subsection*{4) ATTACK PLAN (ROUND-2)}
Round-1 identified the obstruction qualitatively: for $k\ge 3$ the early terms of $R_{p-1,k}$ are too large for the $k=1,2$ argument to work directly.

The exact claims to prove now are:

\begin{itemize}
\item a lower-bound statement showing that truncating after $L+1$ terms with $L\le k-2$ can \emph{never} produce a subunit error term;
\item an upper-bound statement showing that truncating after the first $k$ terms \emph{always} leaves an error of size $O_k(1/p)$.
\end{itemize}

This turns the vague Round-1 obstruction into a sharp theorem: the first unresolved term count is exactly $k$.

\subsection*{5) WORK (ROUND-2)}
Fix $k\ge 2$ and let $p$ be a prime.

\paragraph{Claim 1 (lower bound on each local term).}
For every $j$ with $0\le j\le k-1$,
\[
U_{j,k}(p)\ge p^{k-j-1}.
\]

\paragraph{Proof.}
Since $p+j$ is a divisor of itself,
\[
\sigma_k(p+j)=\sum_{d\mid (p+j)} d^k \ge (p+j)^k.
\]
Also,
\[
p(p+1)\cdots(p+j)\le p\,(p+j)^j.
\]
Therefore
\[
U_{j,k}(p)
=\frac{\sigma_k(p+j)}{p(p+1)\cdots(p+j)}
\ge \frac{(p+j)^k}{p\,(p+j)^j}
=\frac{(p+j)^{k-j}}{p}
\ge p^{k-j-1}.
\]
\hfill$\square$

\medskip

\paragraph{Corollary 2 (fewer than $k$ terms can never give a small remainder).}
Let $L$ be an integer with $0\le L\le k-2$. Then
\[
\sum_{j=L+1}^{\infty}U_{j,k}(p)\ge p^{k-L-2}.
\]
In particular, for every sufficiently large prime $p$ the remainder after truncating after the first $L+1$ terms is larger than $1$.

\paragraph{Proof.}
The omitted tail is positive, so it is at least its first term:
\[
\sum_{j=L+1}^{\infty}U_{j,k}(p)\ge U_{L+1,k}(p).
\]
Apply Claim 1 with $j=L+1$.
\hfill$\square$

This proves a genuine barrier: the Round-1 style argument cannot work after keeping fewer than $k$ local terms.

\medskip

\paragraph{Claim 3 (uniform $O_k(1/p)$ tail after $k$ terms).}
Define
\[
Q_{p,k}:=\sum_{j=0}^{k-1}U_{j,k}(p),
\qquad
E_{p,k}:=\sum_{j=k}^{\infty}U_{j,k}(p).
\]
Then for every prime $p\ge 2^k$,
\[
0<E_{p,k}\le \frac{2^{k+1}\zeta(k)}{p}.
\]

\paragraph{Proof.}
For $k\ge 2$ one has the standard bound
\[
\sigma_k(n)=n^k\sum_{d\mid n}\frac1{(n/d)^k}
\le n^k\sum_{m=1}^{\infty}\frac1{m^k}
=\zeta(k)n^k.
\]
Hence, with
\[
\nu_{j,k}(p):=\zeta(k)\frac{(p+j)^k}{p(p+1)\cdots(p+j)},
\]
we have $U_{j,k}(p)\le \nu_{j,k}(p)$ for every $j$.

First,
\[
\nu_{k,k}(p)
\le \zeta(k)\frac{(p+k)^k}{p^{k+1}}
\le \zeta(k)\frac{(2p)^k}{p^{k+1}}
=\frac{2^k\zeta(k)}{p},
\]
because $p\ge k$.

Second, for every $j\ge k$,
\[
\frac{\nu_{j+1,k}(p)}{\nu_{j,k}(p)}
=\frac{(p+j+1)^k}{(p+j)^k}\cdot\frac1{p+j+1}
=\frac{(p+j+1)^{k-1}}{(p+j)^k}
\le \frac{2^{k-1}}{p+j}
\le \frac12,
\]
since $p\ge 2^k$.

Therefore the majorant tail is geometric:
\[
E_{p,k}
\le \sum_{j=k}^{\infty}\nu_{j,k}(p)
\le 2\nu_{k,k}(p)
\le \frac{2^{k+1}\zeta(k)}{p}.
\]
The positivity $E_{p,k}>0$ is obvious.
\hfill$\square$

\medskip

\paragraph{Corollary 4 (finite-dimensional reduction).}
Assume $\alpha_k\in\mathbf{Q}$. Then for every sufficiently large prime $p$,
\[
R_{p-1,k}\in\mathbf{Z}
\]
and hence
\[
\min_{m\in\mathbf{Z}}
\left|
\sum_{j=0}^{k-1}\frac{\sigma_k(p+j)}{p(p+1)\cdots(p+j)}-m
\right|
\le \frac{2^{k+1}\zeta(k)}{p}.
\]

\paragraph{Proof.}
The Round-1 factorial-tail lemma gives $R_{p-1,k}\in\mathbf{Z}$ for every sufficiently large prime $p$. Since
\[
R_{p-1,k}=Q_{p,k}+E_{p,k},
\]
Claim 3 shows that the distance from $Q_{p,k}$ to the nearest integer is at most $E_{p,k}\le 2^{k+1}\zeta(k)/p$.
\hfill$\square$

\medskip

\paragraph{Interpretation.}
The open part of Problem \#252 is now reduced to the following explicit task:

\begin{quote}
For fixed $k\ge 3$, show that for infinitely many large primes $p$ the finite sum
\[
Q_{p,k}=\sum_{j=0}^{k-1}\frac{\sigma_k(p+j)}{p(p+1)\cdots(p+j)}
\]
is \emph{not} within $O_k(1/p)$ of any integer.
\end{quote}

This is a strict strengthening of the Round-1 obstruction, because it identifies the exact number of local terms that must be controlled.

\subsection*{6) ADVERSARIAL VERIFICATION}

\begin{itemize}
\item Claim 1 only uses the trivial divisor $d=p+j$ and the elementary denominator bound
\[
p(p+1)\cdots(p+j)\le p(p+j)^j.
\]
No hidden hypothesis is present.
\item Corollary 2 is sharp in the sense relevant here: when $L=k-2$, the omitted tail is still at least $1$, so no interval of length $<1$ can arise from such a truncation.
\item Claim 3 uses only the convergent Euler sum $\zeta(k)$ and a geometric-ratio estimate valid for every $p\ge 2^k$.
\item The reduction in Corollary 4 is fully rigorous and depends only on the Round-1 integer-tail lemma.
\end{itemize}

Thus there is no logical gap in the new reduction: the unresolved issue is purely the arithmetic behavior of the explicit $k$-term local sum $Q_{p,k}$.

\subsection*{7) FINAL}
\textbf{UNRESOLVED (BUT STRICTLY ADVANCED).}

The strict Round-2 advance is an exact threshold theorem for the prime-tail method.

\begin{itemize}
\item Truncating after fewer than $k$ local terms can never leave a subunit remainder.
\item Truncating after exactly the first $k$ local terms always leaves an error $O_k(1/p)$.
\item Therefore rationality of $\alpha_k$ would force the explicit $k$-term quantity
\[
Q_{p,k}=\sum_{j=0}^{k-1}\frac{\sigma_k(p+j)}{p(p+1)\cdots(p+j)}
\]
to be $O_k(1/p)$-close to an integer for every sufficiently large prime $p$.
\end{itemize}

This is a genuine finite-dimensional reduction of the unresolved cases.

\subsection*{8) COMPLETION ESTIMATE (MANDATORY)}
COMPLETION: 55\%

\subsection*{9) REFERENCES}
Only references actually used or checked.

\begin{itemize}
\item T. F. Bloom, \emph{Erd\H{o}s Problem \#252}, \texttt{https://www.erdosproblems.com/252}, accessed 2026-03-07.
\item P. Erd\H{o}s, \emph{Problem 4493}, Amer. Math. Monthly 59 (1952), 557--558.
\item J.-C. Schlage-Puchta, \emph{The irrationality of a number theoretical series}, Ramanujan J. 15 (2008), 149--158; arXiv version checked at \texttt{https://arxiv.org/abs/1105.1452}.
\item J. B. Friedlander, F. Luca, and M. Stoiciu, \emph{On the irrationality of a divisor function series}, Integers 7 (2007), A36.
\item K. Pratt, \emph{The irrationality of a divisor function series of Erd\H{o}s and Kac}, arXiv:2209.11124.
\end{itemize}

\bigskip
\hrule
\bigskip

